\section{Conclusion}
\label{sec:conclusion}

We introduced \PTA{} as a deliberately minimal experimental probe for active context acquisition under hidden material dynamics, together with a reproducible cross-material edge-push benchmark in Genesis spanning sand, snow, and elastoplastic media (55pp of scripted-transfer spread). All learned policies are trained on sand only and evaluated zero-shot on the four other splits.

The empirical finding is an asymmetric sign pattern. On elastoplastic (recoverable rebound), \PTA{} delivers a 14.7pp transfer gain, a 14.7pp spill reduction, and a 16.7pp success-rate gain over a matched reactive PPO baseline in the original main protocol; a corrected matched-encoder six-seed audit further yields 93.5\% mean M7 transfer versus 39.4\% for six reactive seeds, with four wins and two near ties. Matched ablations show that removing the probe falls below the reactive baseline, while dropping only the encoder output returns to roughly baseline transfer and remains below full \PTA{}. On ID sand, snow, and two sand parameter shifts, the same architecture underperforms by 13--22pp. We organize this pattern through a \emph{recoverable-deformation hypothesis}: probing appears beneficial when probe-induced perturbations relax on the task timescale and detrimental when they leave persistent disruption. A single-seed privileged-parameter diagnostic that scores 0\% on elastoplastic adds a complementary observation: in our sand-only-trained setup, an in-episode probing trace was a more useful conditioning signal than a passive parameter vector.

\PTA{} is therefore best read as a phenomenon report and diagnostic instrument, not as a deployment-ready manipulation stack. The recoverability hypothesis makes concrete, falsifiable predictions: matched recoverable/nonrecoverable materials should flip the probe sign with measured post-probe relaxation rather than material label; increasing probe amplitude or reducing settle time should turn an elastoplastic gain negative past a residual-deformation threshold; and equal-capacity dummy probes or non-sand training sources should not create a probe advantage on persistent-perturbation regimes. Failure of these directional tests would reject recoverability as the organizing explanation. Code, configurations, and seeds will be released upon acceptance.
